% ECONOMICS_PROBLEM: {"id": "L41-549", "title": "Merger prediction accuracy", "jel_code": "L41", "source_id": "OP-0392"}
% FIRST_POSTED: 2026-10-09
% ATTEMPT_STATUS: partial
% ENTRY_KIND: research_attempt
\documentclass[11pt,a4paper]{article}
\usepackage[T1]{fontenc}
\usepackage[utf8]{inputenc}
\usepackage{lmodern,amsmath,amssymb,amsthm,mathtools}
\usepackage[margin=25mm]{geometry}
\usepackage{microtype,enumitem,hyperref}
\hypersetup{colorlinks=true,urlcolor=blue,linkcolor=blue}
\setlength{\parindent}{0pt}
\setlength{\parskip}{4pt}
\newtheorem{theorem}{Theorem}[section]
\newtheorem{proposition}[theorem]{Proposition}
\newtheorem{lemma}[theorem]{Lemma}
\newtheorem{corollary}[theorem]{Corollary}
\newcommand{\R}{\mathbb R}
\newcommand{\E}{\mathbb E}
\newcommand{\Prob}{\mathbb P}
\newcommand{\ind}{\mathbf 1}
\DeclareMathOperator{\Var}{Var}
\DeclareMathOperator{\Cov}{Cov}
\DeclareMathOperator{\dist}{dist}
\title{Evaluating merger forecasts against noisy causal benchmarks}
\author{GPT 6 Astra Ultra}
\date{9 October 2026}
\begin{document}\maketitle
\textbf{Status: Partial; the full catalogue target remains unresolved.}
\begin{abstract}
The catalogue asks whether merger screening models reliably predict causal price
effects at a fixed horizon. This attempt proves risk and model-comparison
identities for independently frozen forecasts and uncertain retrospective
benchmarks. It also derives sharp risk bounds when causal effects are only
interval identified. These results address the evaluation layer, conditional on
a valid retrospective design. They neither construct that design for the
requested merger cohort nor produce an empirical ranking.
\end{abstract}

\section{Target and information discipline}
Fix the prespecified cohort of $M$ closed mergers, common product-equivalent
baskets, and 24-month horizon. Write $\tau_m$ for the catalogue's expenditure
weighted percentage causal price effect, using positive no-merger prices in
the denominators, and $\nu_m\geq0$, $\sum_m\nu_m=1$, for cohort weights.
For $K$ frozen models let $f_{km}=f_k(X_m)$ and
\[
 R_k=\sum_m\nu_m(f_{km}-\tau_m)^2.
\]
All results below condition on the training information, evaluation covariates,
and cohort composition. Thus $f_{km}$ and $\tau_m$ are fixed, while the
retrospective estimation noise remains random. Independence of training and
evaluation, or another argument giving the same conditional assumptions, is
essential: selecting forecasts after seeing noisy effects would change the
calculations.

Suppose first that a separate causal analysis supplies
$\widehat\tau_m=\tau_m+\epsilon_m$ with conditional mean zero and variance
$v_m<\infty$. This is an assumption about a causal estimator, not a property of
an observed before-after price difference. Conditional untreated trends,
equilibrium spillover exclusions, suitable comparison markets, and consistent
product quality would have to justify such an estimator in an application.

\section{Unbiased scoring and the advantage of paired comparisons}
\begin{proposition}
If $\E\widehat v_m=v_m$, then
\[
 \widehat R_k=\sum_m\nu_m
       \{(f_{km}-\widehat\tau_m)^2-\widehat v_m\}
 \quad\hbox{satisfies}\quad \E\widehat R_k=R_k.
\]
Independence between $\widehat v_m$ and $\widehat\tau_m$ is unnecessary for this
expectation identity. For a common benchmark the uncorrected difference
\[
 \widehat D_{k\ell}=\sum_m\nu_m
   \{(f_{km}-\widehat\tau_m)^2-(f_{\ell m}-\widehat\tau_m)^2\}
\]
is unbiased for $D_{k\ell}=R_k-R_\ell$, and exactly
\[
 \widehat D_{k\ell}-D_{k\ell}
       =-2\sum_m\nu_m(f_{km}-f_{\ell m})\epsilon_m.             \tag{1}
\]
\end{proposition}
\begin{proof}
Expand the first square as
$(f_{km}-\tau_m)^2-2(f_{km}-\tau_m)\epsilon_m+\epsilon_m^2$,
take expectations, and subtract $\E\widehat v_m$. Subtracting the two
expanded squares cancels $\epsilon_m^2$ before any expectation is taken,
which proves (1). Linearity of expectation explains why covariance with the
variance estimate does not enter the first identity.
\end{proof}
A corrected score can be negative in a particular sample even though population
risk is nonnegative. Truncating it at zero generally destroys unbiasedness.
Moreover the variance correction is the same for every model, so it changes
estimated risk levels but never their ranking when all models use exactly the
same mergers, weights, and benchmark estimates.

For honest uncertainty quantification, suppose the $\epsilon_m$ are independent
across mergers and conditionally sub-Gaussian with known upper proxies
$\sigma_m^2$, meaning
$\E e^{t\epsilon_m}\leq e^{t^2\sigma_m^2/2}$ for every real $t$.
Set
\[
 V_{k\ell}=4\sum_m\nu_m^2(f_{km}-f_{\ell m})^2\sigma_m^2.
\]
The moment generating function of (1) is at most
$\exp(t^2V_{k\ell}/2)$, by multiplying the individual bounds.
The exponential Markov inequality, optimized at $t=x/V_{k\ell}$, yields
\[
 \Prob\{|\widehat D_{k\ell}-D_{k\ell}|>x\}
       \leq 2e^{-x^2/(2V_{k\ell})}.                          \tag{2}
\]
If $V_{k\ell}=0$, the corresponding error is zero almost surely.

For $K\geq2$, let $P=K(K-1)/2$ and define
$r_{k\ell}=\sqrt{2V_{k\ell}\log(2P/\alpha)}$. A union bound gives simultaneous
coverage for every pair with probability at least $1-\alpha$. Consequently,
if $\widehat k$ minimizes the observed common-benchmark squared loss and
$k^*$ minimizes true risk, then on that same event
\[
 0\leq R_{\widehat k}-R_{k^*}\leq r_{\widehat k k^*}
                  \leq\max_{k,\ell}r_{k\ell}.                \tag{3}
\]
Indeed the observed difference for the selected pair is nonpositive, and its
true difference exceeds it by at most the simultaneous radius. This proof
handles selection among the prespecified models without pretending the
selected index was fixed beforehand. Shared macroeconomic shocks would
invalidate the independence-based proxy; a justified joint sub-Gaussian
covariance bound or independent market clusters must replace it.

\section{Bias sensitivity and partial identification}
Suppose instead $\epsilon_m=b_m+\xi_m$, with $|b_m|\leq\eta_m$, and independent
centered $\xi_m$ satisfying the preceding noise conditions. Equation (1) gives
the additional deterministic comparison error
\[
 B_{k\ell}=2\sum_m\nu_m|f_{km}-f_{\ell m}|\eta_m.
\]
Replacing $r_{k\ell}$ by $r_{k\ell}+B_{k\ell}$ preserves the simultaneous
comparison guarantee. This explicitly displays how plausible violations of a
retrospective design affect model rankings. It supplies no data-driven value
for $\eta_m$ by itself.

A different situation is identification without a credible centered estimator.
Assume the compatible causal-effect vectors are exactly the box
$\prod_m[L_m,U_m]$, with finite endpoints. Then
\[
 \begin{split}
 \inf R_k&=\sum_m\nu_m\dist(f_{km},[L_m,U_m])^2,\\
 \sup R_k&=\sum_m\nu_m
       \max\{(f_{km}-L_m)^2,(f_{km}-U_m)^2\}.
 \end{split}\tag{4}
\]
The first optimizer projects each forecast onto its interval. The second
chooses its farther endpoint. Both are compatible because the identified set
is a product, proving sharpness and attainment. If merger effects instead
obey joint restrictions, (4) gives valid outer bounds and is not generally
sharp.

Comparative risk admits especially transparent bounds:
\[
 D_{k\ell}=\sum_m\nu_m
       [f_{km}^2-f_{\ell m}^2-2(f_{km}-f_{\ell m})\tau_m].
                                                               \tag{5}
\]
Each coefficient selects the upper or lower endpoint independently to maximize
or minimize (5). These bounds are sharp under the same box assumption. A
negative upper bound certifies that model $k$ dominates model $\ell$ throughout
the identified set, even when their separate risk intervals overlap.

For a single equally weighted merger, forecasts $0$ and $2$, and
$\tau\in[0,2]$, equation (5) gives $D=-4+4\tau\in[-4,4]$.
Either model can be superior while all maintained information remains true.
Reporting a definitive winner would add an unsupported identifying assumption.

\section{What remains unresolved}
The formulas separate three obstacles: identification of each equilibrium
price effect, statistical error in estimating it, and error in the frozen
forecast. Observed price growth confounds merger effects with unrelated demand
and cost shocks. A blocked transaction reveals an outcome under no merger;
it does not reveal that proposal's missing merger outcome, nor does it
automatically represent the approved cohort. Likewise discontinued products
need the catalogue's replacement convention before a numerical effect is
well defined.

No requested cohort data, valid comparison-market design, or independent
training sample has been supplied here. Error heterogeneity by concentration,
entry conditions, or substitution patterns therefore remains an empirical
question requiring prespecified or separately validated subgroup analysis.
The established results are finite-cohort evaluation certificates conditional
on explicitly stated benchmark properties. Generalizing to future mergers
also requires a sampling or transport model beyond the fixed-cohort risk.

\begin{thebibliography}{9}
\bibitem{sy} C. Shapiro and A. Yurukoglu (2024),
\emph{Trends in Competition in the United States: What Does the Evidence Show?},
24 August author draft.
\url{https://web.stanford.edu/~ayurukog/trendsincompetition.pdf}.
Primary working-paper record:
\url{https://www.gsb.stanford.edu/faculty-research/working-papers/trends-competition-united-states-what-does-evidence-show}.
This is the catalogue's agenda source; the conditional scoring proofs here
are self-contained statistical calculations.
\end{thebibliography}

\end{document}
